% --- Archon physics formalization source begin ---
% archon:physics
% archon:covers PhyXMiniProblems/problem_phyx_mini_0457.lean
% archon:source-report reports/phyx_mini/problem_phyx_mini_0457.source.json

\chapter{Physics problem phyx\_mini\_0457}
\label{ch:PhyXMiniProblems_problem_phyx_mini_0457}

\paragraph{Problem source.}
\#\# Physical scenario

A \$10 \textbackslash{}, \textbackslash{}text\{m\}\$-high cylinder, with a cross-sectional area of \$0.1 \textbackslash{}, \textbackslash{}text\{m\}\textasciicircum{}2\$, has a massless piston at the bottom with water at \$20 \textbackslash{}, \textasciicircum{}\{\textbackslash{}circ\}\textbackslash{}text\{C\}\$ on top of it, as shown in figure. Air at \$300 \textbackslash{}, \textbackslash{}text\{K\}\$, with a volume of \$0.3 \textbackslash{}, \textbackslash{}text\{m\}\textasciicircum{}3\$, under the piston is heated so that the piston moves up, spilling the water out over the side.

\#\# Question

Find the total heat transfer to the air when all the water has been pushed out.

\#\# Answer choices

- A: A: 101.3 kJ

- B: B: 433.1 kJ

- C: C: 169.8 kJ

- D: D: 220.7 kJ

\#\# Figure caption (auxiliary; use the image as primary evidence)

The image depicts a vertical cylindrical container with two main sections. The upper section contains a fluid, labeled as "H₂O," representing water. The pressure at the top of this section is indicated as "P₀." Below the water is a horizontal barrier or interface with a series of lines, which separates the water from the section underneath.

The lower section contains "Air," as labeled. The entire setup is enclosed within solid walls, depicted in gray. To the right of the container, there's a downward arrow labeled "g," representing gravitational force. The overall image illustrates a fluid dynamics scenario, likely focusing on the behavior of fluids under gravity and pressure.

\#\# Dataset metadata

Laws of Thermodynamics; Multi-Formula Reasoning, Physical Model Grounding Reasoning

\paragraph{Recorded answer/context.}
D: D: 220.7 kJ

\paragraph{Figure/image path.}
/root/proposal\_for\_physic/science-mango/phyx\_mini\_run/phyx\_data/test\_image/457.png

\paragraph{Formalization target.}
create a compiling Lean file with sorry bodies at `PhyXMiniProblems/problem\_phyx\_mini\_0457.lean`. The Lean declarations must preserve the physical quantities, dimensions or dimensional roles, figure labels, governing-law hypotheses, and final relation expressed by this problem.
Use Mathlib/Physlib names found through LeanExplore where available. If a physics API is missing, introduce faithful local abstractions rather than scalar placeholder aliases.

\begin{theorem}[Physics formalization target]
\label{thm:physics:phyx_mini_0457:target}
\lean{PhyXMiniProblems.ProblemPhyXMini0457.problem_phyx_mini_0457}
\uses{def:physics:phyx-mini-0457:phyxminiproblems-problemphyxmini0457-energyinjoules, def:physics:phyx-mini-0457:phyxminiproblems-problemphyxmini0457-waterexpulsionsetup, def:physics:phyx-mini-0457:phyxminiproblems-problemphyxmini0457-matcheswaterexpulsionscenario, def:physics:phyx-mini-0457:phyxminiproblems-problemphyxmini0457-matchesproblemreadouts, def:physics:phyx-mini-0457:phyxminiproblems-problemphyxmini0457-matchessuppliedwaterpistonfigure, def:physics:phyx-mini-0457:phyxminiproblems-problemphyxmini0457-usesengineeringreferencedata, def:physics:phyx-mini-0457:phyxminiproblems-problemphyxmini0457-hasphysicalwaterexpulsionparameters, def:physics:phyx-mini-0457:phyxminiproblems-problemphyxmini0457-allwaterhasbeenpushedout, def:physics:phyx-mini-0457:phyxminiproblems-problemphyxmini0457-satisfieswaterexpulsionlaws, def:physics:phyx-mini-0457:phyxminiproblems-problemphyxmini0457-isuniquereportedheatchoice}
The assigned autoformalize agent should translate this physics problem into Lean declarations in the covered file, with theorem and lemma proof bodies written as `by sorry`.
\end{theorem}
\begin{proof}
This is an autoformalization task, not a proof task. Produce statements that can later be proved by the physics prover without weakening the physical model.
\end{proof}

% --- Archon declaration topology begin ---
\section{Declaration topology}
\begin{definition}[Length Quantity]
  \label{def:physics:phyx-mini-0457:phyxminiproblems-problemphyxmini0457-lengthquantity}
  \lean{PhyXMiniProblems.ProblemPhyXMini0457.LengthQuantity}
  A nonnegative physical length, independent of a choice of units
\end{definition}
\begin{proof}
  This is the declared model definition of Length Quantity.
\end{proof}
\begin{definition}[Volume Quantity]
  \label{def:physics:phyx-mini-0457:phyxminiproblems-problemphyxmini0457-volumequantity}
  \lean{PhyXMiniProblems.ProblemPhyXMini0457.VolumeQuantity}
  A nonnegative physical volume, with dimension L³
\end{definition}
\begin{proof}
  This is the declared model definition of Volume Quantity.
\end{proof}
\begin{definition}[Mass Quantity]
  \label{def:physics:phyx-mini-0457:phyxminiproblems-problemphyxmini0457-massquantity}
  \lean{PhyXMiniProblems.ProblemPhyXMini0457.MassQuantity}
  A nonnegative physical mass
\end{definition}
\begin{proof}
  This is the declared model definition of Mass Quantity.
\end{proof}
\begin{definition}[Density Quantity]
  \label{def:physics:phyx-mini-0457:phyxminiproblems-problemphyxmini0457-densityquantity}
  \lean{PhyXMiniProblems.ProblemPhyXMini0457.DensityQuantity}
  A nonnegative mass density, with dimension M L⁻³
\end{definition}
\begin{proof}
  This is the declared model definition of Density Quantity.
\end{proof}
\begin{definition}[Acceleration Quantity]
  \label{def:physics:phyx-mini-0457:phyxminiproblems-problemphyxmini0457-accelerationquantity}
  \lean{PhyXMiniProblems.ProblemPhyXMini0457.AccelerationQuantity}
  A nonnegative acceleration magnitude, with dimension L T⁻²
\end{definition}
\begin{proof}
  This is the declared model definition of Acceleration Quantity.
\end{proof}
\begin{definition}[length In Meters]
  \label{def:physics:phyx-mini-0457:phyxminiproblems-problemphyxmini0457-lengthinmeters}
  \lean{PhyXMiniProblems.ProblemPhyXMini0457.lengthInMeters}
  \uses{def:physics:phyx-mini-0457:phyxminiproblems-problemphyxmini0457-lengthquantity}
  Read a physical length in metres
\end{definition}
\begin{proof}
  This is the declared model definition of length In Meters.
\end{proof}
\begin{definition}[area In Square Meters]
  \label{def:physics:phyx-mini-0457:phyxminiproblems-problemphyxmini0457-areainsquaremeters}
  \lean{PhyXMiniProblems.ProblemPhyXMini0457.areaInSquareMeters}
  Read a physical area in square metres
\end{definition}
\begin{proof}
  This is the declared model definition of area In Square Meters.
\end{proof}
\begin{definition}[volume In Cubic Meters]
  \label{def:physics:phyx-mini-0457:phyxminiproblems-problemphyxmini0457-volumeincubicmeters}
  \lean{PhyXMiniProblems.ProblemPhyXMini0457.volumeInCubicMeters}
  \uses{def:physics:phyx-mini-0457:phyxminiproblems-problemphyxmini0457-volumequantity}
  Read a physical volume in cubic metres
\end{definition}
\begin{proof}
  This is the declared model definition of volume In Cubic Meters.
\end{proof}
\begin{definition}[mass In Kilograms]
  \label{def:physics:phyx-mini-0457:phyxminiproblems-problemphyxmini0457-massinkilograms}
  \lean{PhyXMiniProblems.ProblemPhyXMini0457.massInKilograms}
  \uses{def:physics:phyx-mini-0457:phyxminiproblems-problemphyxmini0457-massquantity}
  Read a physical mass in kilograms
\end{definition}
\begin{proof}
  This is the declared model definition of mass In Kilograms.
\end{proof}
\begin{definition}[density In Kilograms Per Cubic Meter]
  \label{def:physics:phyx-mini-0457:phyxminiproblems-problemphyxmini0457-densityinkilogramspercubicmeter}
  \lean{PhyXMiniProblems.ProblemPhyXMini0457.densityInKilogramsPerCubicMeter}
  \uses{def:physics:phyx-mini-0457:phyxminiproblems-problemphyxmini0457-densityquantity}
  Read a physical mass density in kilograms per cubic metre
\end{definition}
\begin{proof}
  This is the declared model definition of density In Kilograms Per Cubic Meter.
\end{proof}
\begin{definition}[acceleration In Meters Per Second Squared]
  \label{def:physics:phyx-mini-0457:phyxminiproblems-problemphyxmini0457-accelerationinmeterspersecondsquared}
  \lean{PhyXMiniProblems.ProblemPhyXMini0457.accelerationInMetersPerSecondSquared}
  \uses{def:physics:phyx-mini-0457:phyxminiproblems-problemphyxmini0457-accelerationquantity}
  Read an acceleration in metres per second squared
\end{definition}
\begin{proof}
  This is the declared model definition of acceleration In Meters Per Second Squared.
\end{proof}
\begin{definition}[pressure In Pascals]
  \label{def:physics:phyx-mini-0457:phyxminiproblems-problemphyxmini0457-pressureinpascals}
  \lean{PhyXMiniProblems.ProblemPhyXMini0457.pressureInPascals}
  Read a Physlib pressure in pascals
\end{definition}
\begin{proof}
  This is the declared model definition of pressure In Pascals.
\end{proof}
\begin{definition}[energy In Joules]
  \label{def:physics:phyx-mini-0457:phyxminiproblems-problemphyxmini0457-energyinjoules}
  \lean{PhyXMiniProblems.ProblemPhyXMini0457.energyInJoules}
  Read a Physlib energy in joules
\end{definition}
\begin{proof}
  This is the declared model definition of energy In Joules.
\end{proof}
\begin{definition}[energy In Kilojoules]
  \label{def:physics:phyx-mini-0457:phyxminiproblems-problemphyxmini0457-energyinkilojoules}
  \lean{PhyXMiniProblems.ProblemPhyXMini0457.energyInKilojoules}
  \uses{def:physics:phyx-mini-0457:phyxminiproblems-problemphyxmini0457-energyinjoules}
  Read a Physlib energy in kilojoules
\end{definition}
\begin{proof}
  This is the declared model definition of energy In Kilojoules.
\end{proof}
\begin{definition}[temperature In Kelvin]
  \label{def:physics:phyx-mini-0457:phyxminiproblems-problemphyxmini0457-temperatureinkelvin}
  \lean{PhyXMiniProblems.ProblemPhyXMini0457.temperatureInKelvin}
  Physlib's Temperature stores an absolute temperature in an arbitrary zero-preserving temperature unit. The explicit storage unit makes the kelvin and Celsius readouts unambiguous
\end{definition}
\begin{proof}
  This is the declared model definition of temperature In Kelvin.
\end{proof}
\begin{definition}[celsius Zero In Kelvin]
  \label{def:physics:phyx-mini-0457:phyxminiproblems-problemphyxmini0457-celsiuszeroinkelvin}
  \lean{PhyXMiniProblems.ProblemPhyXMini0457.celsiusZeroInKelvin}
  The exact offset between zero degrees Celsius and zero kelvin
\end{definition}
\begin{proof}
  This is the declared model definition of celsius Zero In Kelvin.
\end{proof}
\begin{definition}[temperature In Degrees Celsius]
  \label{def:physics:phyx-mini-0457:phyxminiproblems-problemphyxmini0457-temperatureindegreescelsius}
  \lean{PhyXMiniProblems.ProblemPhyXMini0457.temperatureInDegreesCelsius}
  \uses{def:physics:phyx-mini-0457:phyxminiproblems-problemphyxmini0457-temperatureinkelvin, def:physics:phyx-mini-0457:phyxminiproblems-problemphyxmini0457-celsiuszeroinkelvin}
  Read an absolute temperature on the affine Celsius scale
\end{definition}
\begin{proof}
  This is the declared model definition of temperature In Degrees Celsius.
\end{proof}
\begin{definition}[Process State]
  \label{def:physics:phyx-mini-0457:phyxminiproblems-problemphyxmini0457-processstate}
  \lean{PhyXMiniProblems.ProblemPhyXMini0457.ProcessState}
  The initial state and the state at which all water has been expelled
\end{definition}
\begin{proof}
  This is the declared model definition of Process State.
\end{proof}
\begin{definition}[Gas Substance]
  \label{def:physics:phyx-mini-0457:phyxminiproblems-problemphyxmini0457-gassubstance}
  \lean{PhyXMiniProblems.ProblemPhyXMini0457.GasSubstance}
  Material confined below the piston
\end{definition}
\begin{proof}
  This is the declared model definition of Gas Substance.
\end{proof}
\begin{definition}[Upper Liquid]
  \label{def:physics:phyx-mini-0457:phyxminiproblems-problemphyxmini0457-upperliquid}
  \lean{PhyXMiniProblems.ProblemPhyXMini0457.UpperLiquid}
  Liquid initially resting above the piston
\end{definition}
\begin{proof}
  This is the declared model definition of Upper Liquid.
\end{proof}
\begin{definition}[Vertical Direction]
  \label{def:physics:phyx-mini-0457:phyxminiproblems-problemphyxmini0457-verticaldirection}
  \lean{PhyXMiniProblems.ProblemPhyXMini0457.VerticalDirection}
  Orientation of the cylinder axis and of the gravity arrow
\end{definition}
\begin{proof}
  This is the declared model definition of Vertical Direction.
\end{proof}
\begin{definition}[Process Regime]
  \label{def:physics:phyx-mini-0457:phyxminiproblems-problemphyxmini0457-processregime}
  \lean{PhyXMiniProblems.ProblemPhyXMini0457.ProcessRegime}
  The quasistatic idealization used for the moving piston
\end{definition}
\begin{proof}
  This is the declared model definition of Process Regime.
\end{proof}
\begin{definition}[Figure Object]
  \label{def:physics:phyx-mini-0457:phyxminiproblems-problemphyxmini0457-figureobject}
  \lean{PhyXMiniProblems.ProblemPhyXMini0457.FigureObject}
  Physical objects visible in the supplied raster
\end{definition}
\begin{proof}
  This is the declared model definition of Figure Object.
\end{proof}
\begin{definition}[Figure Label]
  \label{def:physics:phyx-mini-0457:phyxminiproblems-problemphyxmini0457-figurelabel}
  \lean{PhyXMiniProblems.ProblemPhyXMini0457.FigureLabel}
  Literal labels visible in the supplied raster
\end{definition}
\begin{proof}
  This is the declared model definition of Figure Label.
\end{proof}
\begin{definition}[Water Piston Cylinder Figure]
  \label{def:physics:phyx-mini-0457:phyxminiproblems-problemphyxmini0457-waterpistoncylinderfigure}
  \lean{PhyXMiniProblems.ProblemPhyXMini0457.WaterPistonCylinderFigure}
  \uses{def:physics:phyx-mini-0457:phyxminiproblems-problemphyxmini0457-verticaldirection, def:physics:phyx-mini-0457:phyxminiproblems-problemphyxmini0457-figureobject, def:physics:phyx-mini-0457:phyxminiproblems-problemphyxmini0457-figurelabel}
  Qualitative and symbolic information transcribed from the primary image
\end{definition}
\begin{proof}
  This is the declared model definition of Water Piston Cylinder Figure.
\end{proof}
\begin{definition}[Air State]
  \label{def:physics:phyx-mini-0457:phyxminiproblems-problemphyxmini0457-airstate}
  \lean{PhyXMiniProblems.ProblemPhyXMini0457.AirState}
  \uses{def:physics:phyx-mini-0457:phyxminiproblems-problemphyxmini0457-volumequantity}
  One equilibrium state of the closed air sample
\end{definition}
\begin{proof}
  This is the declared model definition of Air State.
\end{proof}
\begin{definition}[Water Expulsion Setup]
  \label{def:physics:phyx-mini-0457:phyxminiproblems-problemphyxmini0457-waterexpulsionsetup}
  \lean{PhyXMiniProblems.ProblemPhyXMini0457.WaterExpulsionSetup}
  \uses{def:physics:phyx-mini-0457:phyxminiproblems-problemphyxmini0457-lengthquantity, def:physics:phyx-mini-0457:phyxminiproblems-problemphyxmini0457-massquantity, def:physics:phyx-mini-0457:phyxminiproblems-problemphyxmini0457-densityquantity, def:physics:phyx-mini-0457:phyxminiproblems-problemphyxmini0457-accelerationquantity, def:physics:phyx-mini-0457:phyxminiproblems-problemphyxmini0457-processstate, def:physics:phyx-mini-0457:phyxminiproblems-problemphyxmini0457-gassubstance, def:physics:phyx-mini-0457:phyxminiproblems-problemphyxmini0457-upperliquid, def:physics:phyx-mini-0457:phyxminiproblems-problemphyxmini0457-processregime, def:physics:phyx-mini-0457:phyxminiproblems-problemphyxmini0457-waterpistoncylinderfigure, def:physics:phyx-mini-0457:phyxminiproblems-problemphyxmini0457-airstate}
  The complete physical setup. The two real-valued path functions are SI readouts of a physical water height and pressure as functions of the SI air volume readout; they do not replace the corresponding physical quantities
\end{definition}
\begin{proof}
  This is the declared model definition of Water Expulsion Setup.
\end{proof}
\begin{definition}[air Temperature In Kelvin]
  \label{def:physics:phyx-mini-0457:phyxminiproblems-problemphyxmini0457-airtemperatureinkelvin}
  \lean{PhyXMiniProblems.ProblemPhyXMini0457.airTemperatureInKelvin}
  \uses{def:physics:phyx-mini-0457:phyxminiproblems-problemphyxmini0457-temperatureinkelvin, def:physics:phyx-mini-0457:phyxminiproblems-problemphyxmini0457-processstate, def:physics:phyx-mini-0457:phyxminiproblems-problemphyxmini0457-waterexpulsionsetup}
  Kelvin readout of the air temperature at a process endpoint
\end{definition}
\begin{proof}
  This is the declared model definition of air Temperature In Kelvin.
\end{proof}
\begin{definition}[water Temperature In Degrees Celsius]
  \label{def:physics:phyx-mini-0457:phyxminiproblems-problemphyxmini0457-watertemperatureindegreescelsius}
  \lean{PhyXMiniProblems.ProblemPhyXMini0457.waterTemperatureInDegreesCelsius}
  \uses{def:physics:phyx-mini-0457:phyxminiproblems-problemphyxmini0457-temperatureindegreescelsius, def:physics:phyx-mini-0457:phyxminiproblems-problemphyxmini0457-waterexpulsionsetup}
  Celsius readout of the water temperature
\end{definition}
\begin{proof}
  This is the declared model definition of water Temperature In Degrees Celsius.
\end{proof}
\begin{definition}[Matches Water Expulsion Scenario]
  \label{def:physics:phyx-mini-0457:phyxminiproblems-problemphyxmini0457-matcheswaterexpulsionscenario}
  \lean{PhyXMiniProblems.ProblemPhyXMini0457.MatchesWaterExpulsionScenario}
  \uses{def:physics:phyx-mini-0457:phyxminiproblems-problemphyxmini0457-massinkilograms, def:physics:phyx-mini-0457:phyxminiproblems-problemphyxmini0457-waterexpulsionsetup}
  Qualitative conditions stated or conventionally idealized in the problem
\end{definition}
\begin{proof}
  This is the declared model definition of Matches Water Expulsion Scenario.
\end{proof}
\begin{definition}[Matches Problem Readouts]
  \label{def:physics:phyx-mini-0457:phyxminiproblems-problemphyxmini0457-matchesproblemreadouts}
  \lean{PhyXMiniProblems.ProblemPhyXMini0457.MatchesProblemReadouts}
  \uses{def:physics:phyx-mini-0457:phyxminiproblems-problemphyxmini0457-lengthinmeters, def:physics:phyx-mini-0457:phyxminiproblems-problemphyxmini0457-areainsquaremeters, def:physics:phyx-mini-0457:phyxminiproblems-problemphyxmini0457-volumeincubicmeters, def:physics:phyx-mini-0457:phyxminiproblems-problemphyxmini0457-waterexpulsionsetup, def:physics:phyx-mini-0457:phyxminiproblems-problemphyxmini0457-airtemperatureinkelvin, def:physics:phyx-mini-0457:phyxminiproblems-problemphyxmini0457-watertemperatureindegreescelsius}
  The four numerical readouts explicitly supplied by the prose. Neither the final temperature nor the requested heat transfer appears here
\end{definition}
\begin{proof}
  This is the declared model definition of Matches Problem Readouts.
\end{proof}
\begin{definition}[Matches Supplied Water Piston Figure]
  \label{def:physics:phyx-mini-0457:phyxminiproblems-problemphyxmini0457-matchessuppliedwaterpistonfigure}
  \lean{PhyXMiniProblems.ProblemPhyXMini0457.MatchesSuppliedWaterPistonFigure}
  \uses{def:physics:phyx-mini-0457:phyxminiproblems-problemphyxmini0457-figureobject, def:physics:phyx-mini-0457:phyxminiproblems-problemphyxmini0457-figurelabel, def:physics:phyx-mini-0457:phyxminiproblems-problemphyxmini0457-waterexpulsionsetup}
  Primary-raster evidence: P₀ is applied at the open upper water surface, water is above the piston, air is below it, and g points downward. The image contains no numerical heat-transfer readout
\end{definition}
\begin{proof}
  This is the declared model definition of Matches Supplied Water Piston Figure.
\end{proof}
\begin{definition}[Uses Engineering Reference Data]
  \label{def:physics:phyx-mini-0457:phyxminiproblems-problemphyxmini0457-usesengineeringreferencedata}
  \lean{PhyXMiniProblems.ProblemPhyXMini0457.UsesEngineeringReferenceData}
  \uses{def:physics:phyx-mini-0457:phyxminiproblems-problemphyxmini0457-densityinkilogramspercubicmeter, def:physics:phyx-mini-0457:phyxminiproblems-problemphyxmini0457-accelerationinmeterspersecondsquared, def:physics:phyx-mini-0457:phyxminiproblems-problemphyxmini0457-pressureinpascals, def:physics:phyx-mini-0457:phyxminiproblems-problemphyxmini0457-waterexpulsionsetup}
  Standard engineering reference values used at the precision of the answer choices: P₀ = 101.3 kPa, liquid-water density 1000 kg/m³, g = 9.8 m/s², dry-air R = 287 J/(kg K), and cᵥ/R = 5/2
\end{definition}
\begin{proof}
  This is the declared model definition of Uses Engineering Reference Data.
\end{proof}
\begin{definition}[Has Physical Water Expulsion Parameters]
  \label{def:physics:phyx-mini-0457:phyxminiproblems-problemphyxmini0457-hasphysicalwaterexpulsionparameters}
  \lean{PhyXMiniProblems.ProblemPhyXMini0457.HasPhysicalWaterExpulsionParameters}
  \uses{def:physics:phyx-mini-0457:phyxminiproblems-problemphyxmini0457-lengthinmeters, def:physics:phyx-mini-0457:phyxminiproblems-problemphyxmini0457-areainsquaremeters, def:physics:phyx-mini-0457:phyxminiproblems-problemphyxmini0457-volumeincubicmeters, def:physics:phyx-mini-0457:phyxminiproblems-problemphyxmini0457-massinkilograms, def:physics:phyx-mini-0457:phyxminiproblems-problemphyxmini0457-densityinkilogramspercubicmeter, def:physics:phyx-mini-0457:phyxminiproblems-problemphyxmini0457-accelerationinmeterspersecondsquared, def:physics:phyx-mini-0457:phyxminiproblems-problemphyxmini0457-pressureinpascals, def:physics:phyx-mini-0457:phyxminiproblems-problemphyxmini0457-waterexpulsionsetup, def:physics:phyx-mini-0457:phyxminiproblems-problemphyxmini0457-airtemperatureinkelvin}
  Positivity and endpoint ordering selecting the physical expansion branch
\end{definition}
\begin{proof}
  This is the declared model definition of Has Physical Water Expulsion Parameters.
\end{proof}
\begin{definition}[All Water Has Been Pushed Out]
  \label{def:physics:phyx-mini-0457:phyxminiproblems-problemphyxmini0457-allwaterhasbeenpushedout}
  \lean{PhyXMiniProblems.ProblemPhyXMini0457.AllWaterHasBeenPushedOut}
  \uses{def:physics:phyx-mini-0457:phyxminiproblems-problemphyxmini0457-volumeincubicmeters, def:physics:phyx-mini-0457:phyxminiproblems-problemphyxmini0457-waterexpulsionsetup}
  The stopping event named in the question. It fixes the final water-column height, not the final air temperature or the heat supplied
\end{definition}
\begin{proof}
  This is the declared model definition of All Water Has Been Pushed Out.
\end{proof}
\begin{definition}[Satisfies Water Expulsion Laws]
  \label{def:physics:phyx-mini-0457:phyxminiproblems-problemphyxmini0457-satisfieswaterexpulsionlaws}
  \lean{PhyXMiniProblems.ProblemPhyXMini0457.SatisfiesWaterExpulsionLaws}
  \uses{def:physics:phyx-mini-0457:phyxminiproblems-problemphyxmini0457-lengthinmeters, def:physics:phyx-mini-0457:phyxminiproblems-problemphyxmini0457-areainsquaremeters, def:physics:phyx-mini-0457:phyxminiproblems-problemphyxmini0457-volumeincubicmeters, def:physics:phyx-mini-0457:phyxminiproblems-problemphyxmini0457-massinkilograms, def:physics:phyx-mini-0457:phyxminiproblems-problemphyxmini0457-densityinkilogramspercubicmeter, def:physics:phyx-mini-0457:phyxminiproblems-problemphyxmini0457-accelerationinmeterspersecondsquared, def:physics:phyx-mini-0457:phyxminiproblems-problemphyxmini0457-pressureinpascals, def:physics:phyx-mini-0457:phyxminiproblems-problemphyxmini0457-energyinjoules, def:physics:phyx-mini-0457:phyxminiproblems-problemphyxmini0457-processstate, def:physics:phyx-mini-0457:phyxminiproblems-problemphyxmini0457-waterexpulsionsetup, def:physics:phyx-mini-0457:phyxminiproblems-problemphyxmini0457-airtemperatureinkelvin}
  Governing geometry, hydrostatics, ideal-gas behavior, boundary work, caloric behavior, and the closed-system first law. The internal-energy relation is the general calorically-perfect ideal-air identity ΔU = (cᵥ/R) Δ(PV). The first law uses the convention that both heat transferred to the air and work done by the air are positive. No field assigns a numerical value to the requested heat transfer or chooses D
\end{definition}
\begin{proof}
  This is the declared model definition of Satisfies Water Expulsion Laws.
\end{proof}
\begin{definition}[Answer Choice]
  \label{def:physics:phyx-mini-0457:phyxminiproblems-problemphyxmini0457-answerchoice}
  \lean{PhyXMiniProblems.ProblemPhyXMini0457.AnswerChoice}
  Labels printed beside the four candidate heat-transfer values
\end{definition}
\begin{proof}
  This is the declared model definition of Answer Choice.
\end{proof}
\begin{definition}[displayed Heat Transfer In Kilojoules]
  \label{def:physics:phyx-mini-0457:phyxminiproblems-problemphyxmini0457-displayedheattransferinkilojoules}
  \lean{PhyXMiniProblems.ProblemPhyXMini0457.displayedHeatTransferInKilojoules}
  \uses{def:physics:phyx-mini-0457:phyxminiproblems-problemphyxmini0457-answerchoice}
  Heat transfer in kilojoules printed beside each answer label
\end{definition}
\begin{proof}
  This is the declared model definition of displayed Heat Transfer In Kilojoules.
\end{proof}
\begin{definition}[recorded Dataset Answer]
  \label{def:physics:phyx-mini-0457:phyxminiproblems-problemphyxmini0457-recordeddatasetanswer}
  \lean{PhyXMiniProblems.ProblemPhyXMini0457.recordedDatasetAnswer}
  \uses{def:physics:phyx-mini-0457:phyxminiproblems-problemphyxmini0457-answerchoice}
  Dataset answer metadata, deliberately not used as a theorem premise
\end{definition}
\begin{proof}
  This is the declared model definition of recorded Dataset Answer.
\end{proof}
\begin{definition}[Is Reported Heat Choice]
  \label{def:physics:phyx-mini-0457:phyxminiproblems-problemphyxmini0457-isreportedheatchoice}
  \lean{PhyXMiniProblems.ProblemPhyXMini0457.IsReportedHeatChoice}
  \uses{def:physics:phyx-mini-0457:phyxminiproblems-problemphyxmini0457-energyinkilojoules, def:physics:phyx-mini-0457:phyxminiproblems-problemphyxmini0457-waterexpulsionsetup, def:physics:phyx-mini-0457:phyxminiproblems-problemphyxmini0457-answerchoice, def:physics:phyx-mini-0457:phyxminiproblems-problemphyxmini0457-displayedheattransferinkilojoules}
  Agreement after reporting to the 0.1 kJ precision of the answer list
\end{definition}
\begin{proof}
  This is the declared model definition of Is Reported Heat Choice.
\end{proof}
\begin{definition}[Is Unique Reported Heat Choice]
  \label{def:physics:phyx-mini-0457:phyxminiproblems-problemphyxmini0457-isuniquereportedheatchoice}
  \lean{PhyXMiniProblems.ProblemPhyXMini0457.IsUniqueReportedHeatChoice}
  \uses{def:physics:phyx-mini-0457:phyxminiproblems-problemphyxmini0457-waterexpulsionsetup, def:physics:phyx-mini-0457:phyxminiproblems-problemphyxmini0457-answerchoice, def:physics:phyx-mini-0457:phyxminiproblems-problemphyxmini0457-isreportedheatchoice}
  A displayed choice is the unique value agreeing at the stated precision
\end{definition}
\begin{proof}
  This is the declared model definition of Is Unique Reported Heat Choice.
\end{proof}
\begin{lemma}[final Air Volume In Cubic Meters]
  \label{lem:physics:phyx-mini-0457:phyxminiproblems-problemphyxmini0457-finalairvolumeincubicmeters}
  \lean{PhyXMiniProblems.ProblemPhyXMini0457.finalAirVolumeInCubicMeters}
  \uses{def:physics:phyx-mini-0457:phyxminiproblems-problemphyxmini0457-volumeincubicmeters, def:physics:phyx-mini-0457:phyxminiproblems-problemphyxmini0457-waterexpulsionsetup, def:physics:phyx-mini-0457:phyxminiproblems-problemphyxmini0457-matchesproblemreadouts, def:physics:phyx-mini-0457:phyxminiproblems-problemphyxmini0457-hasphysicalwaterexpulsionparameters, def:physics:phyx-mini-0457:phyxminiproblems-problemphyxmini0457-allwaterhasbeenpushedout, def:physics:phyx-mini-0457:phyxminiproblems-problemphyxmini0457-satisfieswaterexpulsionlaws}
  With no water remaining, the air fills the 10 m × 0.1 m² cylinder, so its final volume is 1 m³. This is a geometric consequence, not an independent final-volume readout
\end{lemma}
\begin{proof}
  The conclusion follows from the governing relations and figure readouts named in the statement of final Air Volume In Cubic Meters.
\end{proof}
\begin{lemma}[heat Transferred To Air In Joules exact]
  \label{lem:physics:phyx-mini-0457:phyxminiproblems-problemphyxmini0457-heattransferredtoairinjoules-exact}
  \lean{PhyXMiniProblems.ProblemPhyXMini0457.heatTransferredToAirInJoules_exact}
  \uses{def:physics:phyx-mini-0457:phyxminiproblems-problemphyxmini0457-energyinjoules, def:physics:phyx-mini-0457:phyxminiproblems-problemphyxmini0457-waterexpulsionsetup, def:physics:phyx-mini-0457:phyxminiproblems-problemphyxmini0457-matchesproblemreadouts, def:physics:phyx-mini-0457:phyxminiproblems-problemphyxmini0457-usesengineeringreferencedata, def:physics:phyx-mini-0457:phyxminiproblems-problemphyxmini0457-hasphysicalwaterexpulsionparameters, def:physics:phyx-mini-0457:phyxminiproblems-problemphyxmini0457-allwaterhasbeenpushedout, def:physics:phyx-mini-0457:phyxminiproblems-problemphyxmini0457-satisfieswaterexpulsionlaws}
  The hydrostatic pressure path, ideal-air caloric relation, and first law give 220745 J = 220.745 kJ before answer-list reporting
\end{lemma}
\begin{proof}
  The conclusion follows from the governing relations and figure readouts named in the statement of heat Transferred To Air In Joules exact.
\end{proof}
% --- Archon declaration topology end ---
% --- Archon physics formalization source end ---
